\vspace{-2mm}
\section{Introduction}
\label{sec:introduction}

Coding agents are being rapidly adopted across the software engineering industry \citep{cursor}, and providing \agentsmds{} like \texttt{AGENTS.md}, a \texttt{README} specifically targeting agents, has become common practice. With various industry leaders \citep{AGENTSMd2025,UsingClaudeMdFiles} recommending this approach to adapt their agents to specific repositories, \agentsmds{} are now supported by most popular agent frameworks, and included in over 60'000 open-source repositories at the time of writing, as reported by \citet{AGENTSMd2025}. 

These \agentsmds{} typically contain a repository overview and information on relevant developer tooling, aiming to help coding agents to navigate a given repository more efficiently, run build and test commands correctly, adhere to style guides and design patterns, and ultimately to solve tasks to the user's satisfaction more frequently.
To date, despite their widespread adoption, the impact of \agentsmds{} on the coding agent's ability to solve complex software engineering tasks has not been rigorously studied.
This is due to two key challenges: i) because of their recent introduction, \agentsmds{} are not available for instances of prior benchmarks, and ii) popular, well-known repositories, typically used to create such benchmarks, are not representative of most codebases. As a result, a rigorous evaluation of the \agentsmds{} used in practice requires a new, complementary benchmark that contains only issues from less popular repositories with developer-committed \agentsmds{}.

\begin{figure*}[t]
    \centering
    \includegraphics[width=.9\linewidth]{figures/overview.pdf}
    \caption{Overview of our evaluation pipeline. We begin with real-world repositories and tasks derived from past pull requests. For each repository state, we generate three settings: \textcircled{\tiny{1}} If a developer-provided \agentsmd{} exists, we include it in the repository. In \textcircled{\tiny{2}}, we omit the \agentsmd{}. \textcircled{\tiny{3}} We use the coding agent's recommended settings to generate the \agentsmd{}. Then we pass the repository and \agentsmd{} to the coding agent and instruct it to autonomously resolve the task. We finally analyze the trace for behavioral changes and apply the generated patch to check for task resolution success.}
    \label{fig:overview}
\end{figure*}


\paragraph{This work: Benchmarking \agentsmds{}' impact on resolving GitHub issues}
In this work, we investigate the effect of actively used \agentsmds{} on the resolution of real-world coding tasks.
We evaluate agents both in popular and less-known repositories, and, importantly, with \agentsmds{} provided by repository developers.
For this purpose, we construct a novel benchmark (\cref{fig:overview}, left), \planbench{}, comprising Python software engineering tasks, created specifically from real GitHub issues.
The benchmark contains $138$ unique instances, covering both bug-fixing and feature addition tasks across $12$ recent and niche repositories, which all feature developer-written \agentsmds{}. \planbench{} complements \swebench{}, which we leverage for the evaluation of automatically generated \agentsmds{} on popular repositories.
We evaluate coding agents in three settings (\cref{fig:overview}, middle): without any \agentsmd{}, with \agentsmds{} automatically generated using agent-developer recommendations, and with the developer-provided \agentsmd{}.
Our code to generate \planbench{} instances and evaluate coding agents is available \href{https://github.com/eth-sri/agentbench}{here}.

Surprisingly, we observe that \emph{developer-provided files only marginally improve performance} compared to omitting them entirely (an increase of 4\% on average), while \emph{LLM-generated \agentsmds{} have a small negative effect} on agent performance (a decrease of 3\% on average).
These observations are robust across different LLMs and prompts used to generate the \agentsmds{}.
In a more detailed analysis (\cref{fig:overview}, right), we observe that \agentsmds{} lead to increased exploration, testing, and reasoning by coding agents, and, as a result, increase costs by over 20\%.
We therefore suggest omitting LLM-generated \agentsmds{} for the time being, contrary to agent developers' recommendations, and including only minimal requirements (\eg specific tooling to use with this repository).
We hope our evaluation framework will aid agent and model developers to improve the helpfulness of LLM-generated \agentsmds{}.

\paragraph{Key contributions} Our key contributions are:
\begin{enumerate}
    \item \planbench{}, a new curated benchmark for the impact of actively used \agentsmds{} on agents' ability to solve real-world software engineering tasks.
    \item An extensive evaluation of different coding agents and underlying models on \planbench{} and \swebench{}, showing that LLM-generated \agentsmds{} tend to decrease agent performance, across models or prompts used to generate them, while developer-written \agentsmds{} tend to slightly improve it.
    \item A detailed investigation of agent traces, showing that \agentsmds{} lead to more thorough testing and exploration by coding agents.
\end{enumerate}
